\chapter{Incidente 5: [título pendiente]}
\label{chap:incidente5}
\section{Contextualización}
\pendiente{Identificar organización, producto o sector, cronología con fechas verificadas, fecha de divulgación y relación con el periodo admitido.}
\section{Descripción técnica}
\pendiente{Explicar tipo de fallo o incidente, activos y versiones afectados, vector de entrada, condiciones necesarias y cadena causal. Diferenciar explotación confirmada de escenarios posibles. Añadir un esquema propio si aporta claridad y citar la evidencia técnica.}
\section{CVE y evaluación CVSS, si procede}
\pendiente{Registrar identificadores verificados y su correspondencia con el caso. Indicar versión CVSS, vector, score, fuente y fecha. Si no hay score publicado, intentar estimarlo con supuestos y razonamiento por métrica; desarrollar el cálculo en el anexo. Justificar expresamente cuando CVE o CVSS no sean aplicables.}
\section{Impacto técnico y social}
\pendiente{Analizar a quién afecta y cómo: confidencialidad, integridad, disponibilidad, servicios esenciales, personas y organizaciones. Separar daño observado de impacto potencial. Fechar cifras, citar su origen y señalar incertidumbre; CVSS no sustituye este análisis.}
\section{Prevención y mitigación}
\pendiente{Distinguir medidas preventivas, contención, corrección y recuperación. Asociar cada medida al mecanismo del fallo, priorizarla y justificar sus límites y viabilidad. Verificar versiones corregidas y recomendaciones oficiales.}
\section{Justificación de la selección}
\pendiente{Aplicar los criterios del capítulo metodológico y explicar qué aporta este caso al conjunto. Evitar justificarlo únicamente por su repercusión mediática.}
\section{Análisis crítico}
\pendiente{Contrastar fuentes y respuestas de los actores; analizar causas de fondo, decisiones y alternativas. Formular una interpretación propia sustentada en evidencia y reconocer qué no puede concluirse.}
